\chapter{ARIES}
\label{ch:aries}

ARIES (Algorithm for Recovery and Isolation Exploiting Semantics), published by
Mohan, Haderle, Lindsay, Pirahesh, and Schwarz in 1992 \cite{aries}, is the
recovery algorithm behind most disk-based database systems. It supports steal
and no force, fine-grained locking, and partial rollbacks, and it recovers
correctly even if the system crashes again \emph{during} recovery. Its central
idea fits in three words: \textbf{repeat history, then undo}.

\section{The data structures}

\begin{description}[leftmargin=1.2em]
\item[Log records] as in \cref{ch:wal}: LSN, transaction id, prevLSN, and type-specific fields.
\item[pageLSN] on every page: the LSN of the last change applied to it.
\item[Transaction table] for every active transaction: its status and \term{lastLSN}, the LSN of its most recent record.
\item[Dirty page table] for every page dirty in the buffer pool: its \term{recLSN}, the LSN of the first record that made it dirty since it was last read from disk. Any change to the page before recLSN is already on disk.
\item[Compensation log records (CLRs)] written when a change is undone. A CLR describes the undo, and carries \term{undoNextLSN}: the next record of that transaction still to undo.
\end{description}

\section{Checkpoints}

Without checkpoints, recovery would read the log from the beginning of time. A
\term{fuzzy checkpoint} writes the transaction table and the dirty page table
into the log without stopping the system or forcing any pages, and records the
checkpoint's LSN in a well-known place. Recovery then starts from the last
checkpoint: the tables in it tell analysis everything that happened before.

\section{Pass 1: Analysis}

Scan the log forward from the last checkpoint, rebuilding the two tables as of
the crash:

\begin{itemize}
\item On any record of a transaction not in the table, add it. Update its lastLSN.
\item On a commit, mark it committed; on an end record, remove it.
\item On an update or CLR for a page not in the dirty page table, add the page with recLSN = this LSN.
\end{itemize}

At the end, every transaction still in the table without a commit is a
\term{loser}. The \term{redoLSN} is the smallest recLSN in the dirty page table:
no earlier change can be missing from disk.

\section{Pass 2: Redo}

Scan forward from redoLSN, and reapply every update and every CLR, for winners
and losers alike, unless the change is provably already on disk:

\begin{lstlisting}[language=pseudo]
for each update or CLR record r, from redoLSN forward:
  if r.page not in dirty page table:          skip   // page was clean at the crash
  if dirtyPageTable[r.page].recLSN > r.lsn:   skip   // flushed after r
  fetch page
  if page.pageLSN >= r.lsn:                   skip   // r already on the page
  apply r to the page; page.pageLSN = r.lsn
\end{lstlisting}

Redo \textbf{repeats history}: it brings the database to exactly its state at
the crash, uncommitted changes included. That sounds wasteful, but it makes undo
simple: undo now starts from a precise, known state, and can use the same
logical operations that a normal abort uses.

\section{Pass 3: Undo}

Undo the losers, newest change first, across all of them at once:

\begin{lstlisting}[language=pseudo]
toUndo = { lastLSN of each loser }
while toUndo is not empty:
  r = record with the largest LSN in toUndo; remove it
  if r is a CLR:
    if r.undoNextLSN exists: add r.undoNextLSN to toUndo
    else: write an end record for r.txn
  else if r is an update:
    undo r; write CLR(undoNextLSN = r.prevLSN)
    if r.prevLSN exists: add r.prevLSN to toUndo
    else: write an end record for r.txn
\end{lstlisting}

\begin{keyidea}
A CLR is redone but never undone. If the system crashes during undo, the next
recovery redoes the CLRs (history repeats), then follows their undoNextLSN
pointers past everything already undone. Recovery work is therefore bounded no
matter how many times the system crashes during recovery.
\end{keyidea}

\section{A worked example}

Two transactions run. $T_1$ commits; $T_2$ does not. The buffer pool happened
to write page P2 (stealing $T_2$'s uncommitted change), but never wrote P1.

\begin{center}
\begin{tabular}{@{}rlll@{}}
\toprule
LSN & Transaction & Record & prevLSN \\
\midrule
10 & $T_1$ & update P1 & --- \\
20 & $T_2$ & update P2 & --- \\
30 & $T_1$ & commit & 10 \\
40 & $T_2$ & update P1 & 20 \\
\multicolumn{4}{@{}l}{\emph{crash.} On disk: P1 has pageLSN 0; P2 has pageLSN 20.} \\
\bottomrule
\end{tabular}
\end{center}

\paragraph{Analysis.} $T_1$ committed; $T_2$ is a loser with lastLSN 40. The
dirty page table is \{P1: recLSN 10, P2: recLSN 20\}, so redoLSN is 10.

\paragraph{Redo.} LSN 10 on P1: pageLSN 0 $<$ 10, so redo; P1's pageLSN becomes
10. LSN 20 on P2: pageLSN 20 $\ge$ 20, so skip; the stolen write already put it
there. LSN 40 on P1: pageLSN 10 $<$ 40, so redo; pageLSN becomes 40. The
database is now exactly as it was at the crash, including $T_2$'s two changes.

\paragraph{Undo.} toUndo = \{40\}. Undo LSN 40 on P1, write CLR 50 with
undoNextLSN 20. Undo LSN 20 on P2, write CLR 60 with no undoNextLSN, then an end
record for $T_2$. $T_1$'s change survives; $T_2$ left no trace.

If the system had crashed right after CLR 50, the next analysis would find $T_2$
with lastLSN 50; redo would reapply CLR 50; undo would read CLR 50, jump straight
to LSN 20, and never undo LSN 40 a second time.

\section{Structure changes: nested top actions}

A B+ tree split moves keys between pages. Suppose $T$ inserts a key, which
splits a leaf, and later aborts. Undoing ``the split'' would be wrong: another
transaction may have committed an insert into the new leaf since then. The
split is a change to the \emph{structure}, which is valid regardless of who
caused it; only the \emph{logical} insert belongs to $T$.

ARIES logs a split as a \term{nested top action}: the split's page records,
followed by a dummy CLR whose undoNextLSN points to the record just before the
split. Redo replays the split like any other change. Undo, following the dummy
CLR, skips over it. The logical ``insert key $k$'' is logged separately, and its
undo is the logical ``delete key $k$'', performed wherever that key now lives,
which is why undo in ARIES is logical while redo is physical.

\begin{ontodb}[: ARIES]
OntoDB's page-based recovery follows this chapter: steal and no force,
page-oriented redo, logical undo of triples, CLRs that are redone and never
undone, and B+ tree splits logged as redo-only nested top actions with the
triple insert logged separately. On open, recovery prints one line per pass
(\code{analysis}, \code{redo}, \code{undo}). The bar for the feature is a crash
harness that kills the process at random points (including named crash points
inside splits and between a log flush and a page write) and compares the
reopened store with an oracle that replayed only acknowledged commits:
\code{scripts/crashtest --seeds 200} must pass.
\end{ontodb}

\begin{decision}{A blocking checkpoint first (planned)}
\dpart{Choice} The page store's \code{\textbackslash checkpoint} may pause
updates while it writes the checkpoint record; the comment above
\code{Checkpoint} names which kind it is.
\dpart{Alternatives} ARIES's fuzzy checkpoint, which never stops the system.
\dpart{Why} A blocking checkpoint is easy to reason about and enough to bound
recovery. Fuzzy checkpointing is an optimization to add once the three passes
are correct.
\dpart{Cost} Updates stall for the length of the checkpoint.
\end{decision}

\begin{decision}{Recovery is tested against an oracle}
\dpart{Choice} A crash harness kills the process at random points and at named
crash points, reopens the store, and compares it with an oracle that applied
only acknowledged commits. The harness is itself tested on deliberately broken
stores.
\dpart{Alternatives} Hand-written recovery scenarios only.
\dpart{Why} The bugs in recovery live in interleavings nobody thinks to write.
Many random crash points find them; a harness that catches a planted lost
commit proves it can.
\dpart{Cost} Slow tests, and failures that need a seed to reproduce.
\end{decision}

\begin{status}
The log record types used here, including CLRs and checkpoints, exist and
serialize (plan 8.1). Abort by the log, checkpoints, the three passes, and the
crash test are plans 8.4 to 8.7, after the buffer pool and B+ tree. The crash
harness already detects a lost commit and a surviving uncommitted write on
deliberately faulty test stores.
\end{status}

\begin{checkyourself}
\item In the worked example, suppose the pool had also written P1 after LSN 40. What does redo do now, and what does undo do?
\item Why does ARIES redo the losers' changes, only to undo them right after?
\item The system crashes after CLR 60 but before $T_2$'s end record. What does the next recovery do?
\item Why is a B+ tree split never undone, even when the transaction that caused it aborts?
\item What do the dirty page table's recLSNs let redo skip?
\end{checkyourself}

\section*{Further reading}
Mohan et al.\ \cite{aries}, the whole paper; Ramakrishnan and Gehrke
\cite{cowbook}, chapter 18, for a gentler walkthrough; Petrov \cite{petrov},
chapter 5.
